\documentclass[10pt,a4paper]{amsart}
\usepackage[margin=19mm]{geometry}
\usepackage{amsmath,amssymb,mathtools}
\usepackage[scr=rsfs,cal=euler]{mathalfa}
\usepackage{xfrac}
\usepackage[unicode,hidelinks,bookmarks=false]{hyperref}
\newcommand{\ie}{i.e.\ }
\newcommand{\cf}{cf.\ }
\newcommand{\resp}{respectively\ }
\newcommand{\Subsection}{\subsection}
\providecommand{\nobreakdash}{\nobreak\hbox{-}}
\setlength{\emergencystretch}{2em}
\multlinegap0pt
\usepackage{amssymb}
\newtheorem{prop}{Proposition}
\title{Homogeneous Freudenthal algebras and the first Tits construction}
\author{Holger P. Petersson, Maneesh Thakur}
\date{Source: arXiv:2603.17649v1 (CC BY 4.0)}
\begin{document}
\maketitle

\noindent\emph{Source and licence.} Homogeneous Freudenthal algebras and the first Tits construction. Version arXiv:2603.17649v1 (CC BY 4.0), distributed by arXiv under the Creative Commons Attribution 4.0 International licence, http://creativecommons.org/licenses/by/4.0/, as recorded in the arXiv licence metadata for this exact version and shown on the abstract page https://arxiv.org/abs/2603.17649v1. Full licence notice: \texttt{licenses/arxiv-2603.17649-CC-BY-4.0.txt}.\par\smallskip

\noindent\textit{Editorial scope.} Counted unit: the article's prop key (source0.tex lines 1039--1041). The article's complete original proof of it is delivered with it (source0.tex lines 1043--1072, one \texttt{proof} environment of 30 lines). No mathematics is written, repaired or re-derived: the statement and the proof are the article's own lines, in the article's own notation, and no retained line is substituted or re-flowed.  No further source-local context is needed: every cross-reference of every retained line resolves inside the two delivered spans.  Nothing was omitted from any retained span, so the package carries no omission record.  No retained line contains a citation, so the wrapper carries no bibliography and no bibliography entry.  No retained line contains a figure environment, an image-inclusion command or drawing code, so no image is delivered and the package declares no asset dependency.  Editorial formatting only, in addition: an AMS \texttt{amsart} preamble that restates the article's own declarations and macros used by the retained lines, namely the environments \texttt{prop} on separate counters, together with the standard packages those lines need.  The retained environments print in this document as Proposition 1 in order of appearance, whereas the article prints the same items with the numbering of its own full document.  The complete native source of the article as distributed in the arXiv e-print for this exact version is bundled unchanged as \texttt{sources/196636/source0.tex}.
\begin{prop}\label{key}
	Let $q:V \to k$ be an anisotropic quadratic form over $k$ and suppose ${\bf s},{\bf t}$	are independent indeterminates. Then the quadratic form $q_K \oplus {\bf s} q_K$ over the field $K := k({\bf s},{\bf t})$ does not represent the element ${\bf t} \in K$.
\end{prop}

\begin{proof} We always write $q$ instead of $q_K$ for simplicity. Arguing indirectly, there are elements  $x,y \in V_K$ such that $q(x) + {\bf s} q(y) = {\bf t}$. Then $x \ne 0$ or $y \ne 0$.	Writing
	\begin{align*}
		x = \frac{u({\bf s},{\bf t})}{g({\bf s},{\bf t})}, \quad y = \frac{v({\bf s},{\bf t})}{g({\bf s},{\bf t})}
	\end{align*}
	with $u({\bf s},{\bf t}),v({\bf s},{\bf t}) \in V[{\bf s},{\bf t}] := V \otimes_k k[{\bf s},{\bf t}]$ not both zero and $0 \ne g({\bf s},{\bf t}) \in k[{\bf s},{\bf t}]$, we may clear denominators to conclude
	\begin{align}
		\label{ENCE} q\big(u({\bf s},{\bf t})\big) + {\bf s} q\big(v({\bf s},{\bf t})\big) = {\bf t} g({\bf s},{\bf t})^2.
	\end{align}
	For some $m,n,p \in \mathbb{N}$, we have the expansions
	\begin{align*}
		u({\bf s},{\bf t}) =\,\,&u_m({\bf s}){\bf t}^m + \cdots, \\
		v({\bf s},{\bf t}) =\,\,&v_n(s){\bf t}^n + \cdots, \\
		g({\bf s},{\bf t}) =\,\,&g_p({\bf s}){\bf t}^p + \cdots,  
	\end{align*}
	where $\cdots$ refers to lower terms in ${\bf t}$, $u_m({\bf s}), v_n({\bf s}) \in V[{\bf s}]$, $0 \ne g_p({\bf s}) \in k[{\bf s}]$ and $u_m({\bf s}) \ne 0$ (resp. $v_n({\bf s}) \ne 0$) if $u({\bf s},{\bf t}) \ne 0$ (resp. $v({\bf s},{\bf t}) \ne0$). We conclude
	\begin{align*}
		q\big(u({\bf s},{\bf t})\big) =\,\,&q\big(u_m({\bf s})\big){\bf t}^{2m} + \cdots, \\
		q\big(v({\bf s},{\bf t})\big) =\,\,& q\big(v_n({\bf s})\big){\bf t}^{2n} + \cdots, \\
		g({\bf s},{\bf t})^2 =\,\,&g_p({\bf s})^2{\bf t}^{2p} + \cdots,
	\end{align*}
	where $q(u_m({\bf s}))$ (resp. $q(v_n({\bf s}))$) is non-zero whenever $u({\bf s},{\bf t})$ (resp. $v({\bf s},{\bf t})$) is since $q_{k({\bf s})}$ continues to be anisotropic. Plugging this into \eqref{ENCE}, we deduce
	\begin{align}
		\label{HIGHTE} q\big(u_m({\bf s})\big){\bf t}^{2m} + \cdots + {\bf s} q\big(v_n({\bf s})\big){\bf t}^{2n} + \cdots = g_p({\bf s})^2{\bf t}^{2p+1} + \cdots.
	\end{align}
	Here the degree in ${\bf t}$ of the right-hand side is odd, so we will arrive at a contradiction once we have shown that the degree in ${\bf t}$ of the left-hand side is even. This is clear if $m \ne n$. On the other hand, if $m = n$, then the coefficient of the highest term in ${\bf t}$ of the left-hand side is 
	\[
	q\big(u_m({\bf s})\big) + {\bf s} q\big(v_n({\bf s})\big),
	\]
	which cannot be zero unless both summands are non-zero. But then it is the sum of two non-zero polynomials in ${\bf s}$, the first summand being of even, the second summand being of odd degree. Hence the sum itself is not zero, and we have shown in all cases that the left-hand side of \eqref{HIGHTE} has even degree in ${\bf t}$. 
\end{proof}

\end{document}
